% !TEX root = ../../main.tex

\section{Introduction}

The construction of minimal hypersurfaces via min-max methods involves two essentially independent components: a \emph{sweepout framework} that produces a candidate varifold, and a \emph{regularity theory} that promotes this varifold to a smooth hypersurface.

In the classical Almgren--Pitts theory~\cite{Alm65, Pit81}, the sweepout framework operates in the space of integral currents: a combinatorial selection procedure extracts a min-max sequence whose limit satisfies the \emph{almost minimizing} property---the varifold is almost area-minimizing in small annuli. This property feeds directly into the Schoen--Simon regularity theory~\cite{SS81}, which requires either the almost minimizing condition or, more generally, a strong a priori bound on the singular set ($\mathcal{H}^{n-2}(\sing V) < \infty$). Because the construction uses integral currents, integrality of the limit is built in from the outset.

Following the ideas of Simon--Smith~\cite{Smi83}, Colding--de~Lellis~\cite{CL03} and de~Lellis--Tasnady~\cite{DLT13} replaced the discrete setting of Almgren--Pitts with continuous sweepouts by Caccioppoli sets, considerably simplifying the combinatorial machinery. However, they still rely on the same regularity pipeline: an almost minimizing property plus Schoen--Simon.

Chodosh--Liokumovich--Spolaor~\cite{CLS22} introduced a fundamentally different sweepout framework based on \emph{non-excessive sweepouts}---optimal nested sweepouts whose critical slices cannot be locally replaced by competitors of strictly smaller mass. This framework is considerably more streamlined: it works directly with Caccioppoli sets, avoids the combinatorial machinery of Almgren--Pitts, and replaces the almost minimizing property with a more intuitive geometric condition---one-sided homotopic minimization. However, the regularity step in~\cite{CLS22} still invokes the Almgren--Pitts pull-tight and regularity theory as a black box~\cite[Section~3]{CLS22}, leaving the non-excessive framework dependent on the classical machinery it was designed to simplify.

In this paper, we explore the extent to which this dependence can be removed. We construct an alternative regularity pipeline that replaces the Almgren--Pitts combinatorial selection and Schoen--Simon regularity: stationarity follows from the pull-tight argument of Colding--de~Lellis~\cite{CL03}; integrality on the reduced boundary $\partial^*\Omega(x_0)$ of the flat limit follows unconditionally from a lower-semicontinuity comparison $\|V\| \geq |D\chi_{\Omega(x_0)}|$ (Lemma~\ref{lem:density-lower-bound}) combined with a graphical sheet decomposition (Lemma~\ref{lem:sheet-decomposition}); on the \emph{cancelled mass} $\mu_{\mathrm{c}}$ that appears only when mass cancellation occurs, integrality is conditional on positive density $\|V\|$-a.e.\ on the cancellation locus (Proposition~\ref{prop:positive-density}); and regularity from Wickramasekera's theorem~\cite{Wic14} for stable codimension~$1$ integral varifolds.

Wickramasekera's theorem~\cite{Wic14} applies to stable integral varifolds satisfying the \emph{$\alpha$-structural hypothesis} (Definition~\ref{defn:alpha-structural}), without requiring the almost minimizing property or an a priori bound on the singular set. We verify these hypotheses uniformly in both the mass-cancellation and no-mass-cancellation cases (see Section~\ref{sec:regularity-tools} for a detailed comparison with the classical approach). Combining these ingredients, we obtain:

\begin{theorem}[Regularity for non-excessive min-max]\label{thm:non-excessive-main}
    Let $(M^{n+1}, g)$ be a closed Riemannian manifold with $n \geq 2$. Let $\Phi$ be a non-excessive ONVP sweepout, and let $V$ be the varifold limit of a min-max sequence $x_i \to x_0$. Then $V$ is a stationary $n$-varifold with $\|V\|(M) = W$.
    \begin{enumerate}[label=\textup{(\alph*)}]
        \item (\emph{No mass cancellation.}) If $\mathbf{M}(\Phi(x_0)) = W$, then $V$ is integral with $V = |\partial^*\Omega(x_0)|$, and $\sing V = \varnothing$ for $2 \leq n \leq 6$, $\sing V$ is discrete for $n = 7$, and $\mathcal{H}^{n-7+\gamma}(\sing V) = 0$ for every $\gamma > 0$ when $n \geq 8$. In particular, for $2 \leq n \leq 6$, $\spt\|V\|$ is a smooth, closed, embedded minimal hypersurface.
        \item (\emph{Mass cancellation, conditional.}) If $\mathbf{M}(\Phi(x_0)) < W$, then $\|V\| \geq |D\chi_{\Omega(x_0)}|$ pointwise and the restriction of $V$ to $\partial^*\Omega(x_0) \cap \spt\|V\|$ is integral. The same regularity conclusions hold for $V$ provided the cancelled mass $\mu_{\mathrm{c}} := \|V\| - |D\chi_{\Omega(x_0)}|$ is integral; this reduces to Proposition~\ref{prop:positive-density} (positive density of $V$ for $\|V\|$-a.e.\ point) via Allard's rectifiability theorem and the sheet decomposition (Lemma~\ref{lem:sheet-decomposition}). \CHECK{Theorem~\ref{thm:non-excessive-main}(b) currently states that the cancellation case is conditional on Proposition~\ref{prop:positive-density}. Now that Conjecture~5.9 is upgraded to Proposition with skeleton proof, decide whether the conditionality is: (a) on the proof skeleton being completed (current marked GAPs being closed), or (b) keep as conditional and reference forward to the skeleton. For working draft, option (b) is cleaner: Theorem~\ref{thm:non-excessive-main}(b) remains conditional pending the skeleton being completed.} The proposition is parallel to the multiplicity-one question; see Remark~\ref{rem:cancelled-mass-gap}.
    \end{enumerate}
\end{theorem}

\begin{remark}\label{rem:existence-vs-regularity}
The existence of non-excessive ONVP sweepouts is guaranteed by~\cite[Theorem~2.2]{CLS22} (see Theorem~\ref{thm:non-excessive-existence}). Among the three regularity inputs---integrality, the $\alpha$-structural hypothesis, and stability---only stability uses the ONVP structure. Integrality and the $\alpha$-structural verification require only that the sweepout consists of Caccioppoli sets (see Remark~\ref{rem:integrality-independence}). Nestedness also enters the parity analysis underlying the multiplicity-one question.
\end{remark}

The proof proceeds in three steps: integrality of the limit varifold (Theorem~\ref{thm:integrality}); stability from the non-excessive structure via one-sided homotopic minimization (Proposition~\ref{prop:stability}); and verification of the $\alpha$-structural hypothesis by ruling out junctions through a balancing condition and chord-beats-arc area comparisons (Theorem~\ref{thm:alpha-hypo}). Wickramasekera's theorem~\cite{Wic14} then yields the stated regularity on the part of $V$ where integrality is established; in the mass-cancellation case, the same conclusion holds on the reduced-boundary part of $\|V\|$ unconditionally, and on all of $\|V\|$ conditional on Proposition~\ref{prop:positive-density}.

Two parallel open questions remain in the mass-cancellation case. The first is \emph{multiplicity one}: whether the limit varifold $V$ has multiplicity~$1$, equivalently $\mathbf{M}(\Phi(x_0)) = W$ (no mass cancellation). The second is Proposition~\ref{prop:positive-density} (the proof skeleton appears in Section~\ref{subsec:cancelled-mass-open}): whether, under cancellation, the limit varifold has $\Theta(\|V\|, p) > 0$ for $\|V\|$-a.e.\ $p$, which (via Allard's rectifiability theorem and the sheet decomposition) reduces the integrality of $\mu_{\mathrm{c}} = \|V\| - |D\chi_{\Omega(x_0)}|$ to a routine consequence (Remark~\ref{rem:cancelled-mass-gap}). A positive answer to multiplicity one subsumes both questions and renders the conditional part of Theorem~\ref{thm:non-excessive-main} vacuous.

The paper is organized as follows. Section~\ref{sec:p2-prelim} collects the basic definitions from geometric measure theory: Caccioppoli sets, varifolds, density, and tilt-excess. Section~\ref{sec:non-excessive-prelim} sets up the min-max framework of~\cite{CLS22}: sweepouts, width, ONVP sweepouts, the non-excessive condition, the pull-tight argument, and homotopic minimizers as the geometric mechanism through which non-excessiveness controls regularity. Section~\ref{sec:regularity-tools} compares the classical Schoen--Simon regularity theory with Wickramasekera's theorem for stable codimension~$1$ integral varifolds, and explains why the latter is naturally suited to the non-excessive setting. Section~\ref{sec:integrality} establishes integrality of the limit varifold in both the no-mass-cancellation and mass-cancellation cases. Section~\ref{sec:regularity} verifies stability and the $\alpha$-structural hypothesis, and completes the proof of Theorem~\ref{thm:non-excessive-main}: both cases are handled uniformly by Wickramasekera's theorem. We conclude by discussing the remaining obstruction to multiplicity one.
